\section{Jollof Rice}

\subsection*{Ingredients}
\begin{multicols}{2}
\begin{ingredients}
    \item 200g Guanciale
    \item 2 cups Rice
    \item Red pepper
    \item Tomato
    \item Tomato paste
    \item 2 onions
    \item Scotch bonnet pepper
    \item Garlic
    \item Chicken stock
    \item Neutral oil (vegetable/sunflower)
    \item Herbs: thyme, bay leaf
    \item Spices: fennel, clove, ginger, paprika, tumeric
    \item Salt \& peper
\end{ingredients}
\end{multicols}

\subsection*{Recipe}
\begin{enumerate}
    \item Sauce base: blend the red pepper, 1 onion, tomato with some ckicken stock. Fry off and put aside. 
    \item Spice blend: In the same pan, fry the remaining onion, garlic, tomato paste, scotch bonnet and spices until fragrant.
    \item Combine: Add the rice to the spice mix \& coat the grains. Add half of the sauce base, stock and water. Add the bay leaves.
    \item Cook: Bring to the boil, then simmer on low heat, covered, until rice is cooked (15-20 mins).
    \item Finish: Stir through the remaining sauce base, thyme, adjust seasoning and serve.
\end{enumerate}

\subsection*{Recipe Development}
\begin{itemize}
    \item 31/12/25: Less spice, more tomato \& pepper. More sauce base to serve on the side.
\end{itemize}

\subsection*{Hints}
\begin{itemize}
    \item Not too much spice or tumeric - Jollof is meant to be red, not orange!
    \item Excess sauce base can be served on the side.
    \item Add oil or butter towards the end of cooking, to prevent the rice becoming dry
\end{itemize}
\newpage